\documentclass[10pt,a4paper]{amsart}
\usepackage[margin=19mm]{geometry}
\usepackage{amsmath,amssymb,mathtools}
\usepackage[scr=rsfs,cal=euler]{mathalfa}
\usepackage{xfrac}
\usepackage[unicode,hidelinks,bookmarks=false]{hyperref}
\newcommand{\ie}{i.e.\ }
\newcommand{\cf}{cf.\ }
\newcommand{\resp}{respectively\ }
\newcommand{\Subsection}{\subsection}
\providecommand{\nobreakdash}{\nobreak\hbox{-}}
\setlength{\emergencystretch}{2em}
\multlinegap0pt
\usepackage{tikz}
\usetikzlibrary{cd}
\usepackage{amssymb}
\usepackage[shortlabels]{enumitem}
\usepackage{tikz}
\newtheorem{theorem}{Theorem}
\DeclareMathOperator{\depth}{depth}
\DeclareMathOperator{\height}{ht}
\DeclareMathOperator{\jac}{Jac}
\DeclareMathOperator{\spec}{Spec}
\title{Cohen--Macaulay and Gorenstein Properties of Bi-Amalgamated Algebras with Applications to Algebroid Curves}
\author{Efe G\"urel, Abuzer G\"und\"uz}
\date{Source: arXiv:2604.00266v1 (CC BY 4.0)}
\begin{document}
\maketitle

\noindent\emph{Source and licence.} Cohen--Macaulay and Gorenstein Properties of Bi-Amalgamated Algebras with Applications to Algebroid Curves. Version arXiv:2604.00266v1 (CC BY 4.0), distributed by arXiv under the Creative Commons Attribution 4.0 International licence, http://creativecommons.org/licenses/by/4.0/, as recorded in the arXiv licence metadata for this exact version and shown on the abstract page https://arxiv.org/abs/2604.00266v1. Full licence notice: \texttt{licenses/arxiv-2604.00266-CC-BY-4.0.txt}.\par\smallskip

\noindent\textit{Editorial scope.} Counted unit: the article's theorem cohenmacthm (source0.tex lines 350--357). The article's complete original proof of it is delivered with it (source0.tex lines 358--377, one \texttt{proof} environment of 20 lines). No mathematics is written, repaired or re-derived: the statement and the proof are the article's own lines, in the article's own notation, and no retained line is substituted or re-flowed.  Retained uncounted as the source-local context the proof needs: lines 209--217; lines 259--284.  Nothing was omitted from any retained span, so the package carries no omission record.  No retained line contains a citation, so the wrapper carries no bibliography and no bibliography entry.  The retained lines contain the article's own diagram code, which the wrapper typesets with the standard tikz packages; no image file is delivered and the package declares no asset dependency.  Editorial formatting only, in addition: an AMS \texttt{amsart} preamble that restates the article's own declarations and macros used by the retained lines, namely the environments \texttt{theorem} on separate counters, together with the standard packages those lines need.  The retained environments print in this document as Theorem 1 in order of appearance, whereas the article prints the same items with the numbering of its own full document.  The complete native source of the article as distributed in the arXiv e-print for this exact version is bundled unchanged as \texttt{sources/197581/source0.tex}.
\begin{equation}\label{exactseq}
    \begin{tikzcd}
	0 & {J\oplus J'} & {A\bowtie^{f,g}(J,J')} & {A/I_0} & 0
	\arrow[from=1-1, to=1-2]
	\arrow["i", from=1-2, to=1-3]
	\arrow["\pi", from=1-3, to=1-4]
	\arrow[from=1-4, to=1-5]
\end{tikzcd}
\end{equation}

\begin{theorem}\label{dimdepthcalclemma}
    Let $(A,\mathfrak{m})$ be a local ring with $J\subseteq \jac(B)$ and $J'\subseteq \jac(C)$. Then we have the following.
    \begin{enumerate}
        \item Let $J$ and $J'$ be finitely generated $A$--modules. Then 
        \begin{align*}
            \dim A\bowtie^{f,g}(J,J')=\dim A/(\ker f\cap \ker g).
        \end{align*}
        Furthermore,
        \begin{align*}
            \depth A\bowtie^{f,g}(J,J')=\min(\depth A/I_0,\depth (J\oplus J'))
        \end{align*}
        except possibly when $\depth A/I_0+1=\depth (J\oplus J')$ and the connecting homomorphism 
    \begin{align*}
        \partial^{\depth A/I_0}:H_{\mathfrak{m}}^{\depth A/I_0}(A/I_0)\to H_{\mathfrak{m}}^{\depth A/I_0+1}(J\oplus J')
    \end{align*}
    obtained by applying the functor $H_{\mathfrak{m}}$ to (\ref{exactseq}) has a trivial kernel, in which case $\depth A\bowtie^{f,g}(J,J')=\depth A/I_0+1$.
        \item Let $I_0= \ker f\cap\ker g$. Assume that $f^{-1}(\mathfrak{q}) \neq \mathfrak{m}$ and $g^{-1}(\mathfrak{q}') \neq \mathfrak{m}$ for all $\mathfrak{q}\in \spec (B)\setminus V(J)$ and $\mathfrak{q}'\in\spec (C)\setminus V(J')$, respectively. Then 
        \begin{align*}
            \dim A\bowtie^{f,g}(J,J')= \dim A/(\ker f\cap \ker g)
        \end{align*}
        and
        \begin{align*}
            \depth A\bowtie^{f,g}(J,J')=\min( \depth A/(\ker f\cap \ker g),\depth (J\oplus J')).
        \end{align*}
    \end{enumerate}
\end{theorem}

\begin{theorem}\label{cohenmacthm}
    Let $(A,\mathfrak{m})$ be a local ring with $J\subseteq \jac(B)$ and $J'\subseteq \jac(C)$. Then we have the following.
    \begin{enumerate}
        \item Let $J$ and $J'$ be finitely generated $A$--modules. Assume $\height(\ker f\cap \ker g)=0$. Then, $A\bowtie^{f,g}(J,J')$ is Cohen--Macaulay if and only if $A$ is Cohen--Macaulay, $J$ and $J'$ are maximal Cohen--Macaulay modules with $\depth A/I_0=\depth A$ in the usual case and $\depth A/I_0=\depth A-1$ in the exceptional case (where the exceptional case refers to the condition detailed in Theorem \ref{dimdepthcalclemma}(1)).

        \item  Let $I_0= \ker f\cap\ker g$. Assume that $f^{-1}(\mathfrak{q}) \neq \mathfrak{m}$ and $g^{-1}(\mathfrak{q}') \neq \mathfrak{m}$ for all $\mathfrak{q}\in \spec (B)\setminus V(J)$ and $\mathfrak{q}'\in\spec (C)\setminus V(J')$, respectively. Then, if $A\bowtie^{f,g}(J,J')$ is Cohen--Macaulay, it follows that $A/(\ker f\cap \ker g)$ is Cohen--Macaulay. If furthermore $\depth J,\depth J'<\infty$, then $A\bowtie^{f,g}(J,J')$ is Cohen--Macaulay if and only if $A/(\ker f\cap \ker g)$ is Cohen--Macaulay, and $J$ and $J'$ are big Cohen--Macaulay $A/(\ker f\cap \ker g)$--modules. 
    \end{enumerate}
\end{theorem}

\begin{proof}
     (1) In this case, since $A$ is Cohen--Macaulay and $\height(\ker f \cap \ker g) = 0$, we have $\dim A/(\ker f \cap \ker g) = \dim A$. Thus, by Theorem \ref{dimdepthcalclemma}, $\dim A\bowtie^{f,g}(J,J')=\dim A$. The exceptional case is trivial, so let us assume the usual case. Assume that $A\bowtie^{f,g}(J,J')$ is Cohen--Macaulay, so we must have $\depth A\bowtie^{f,g}(J,J')=\dim A$. First notice that
        \begin{align*}
            \depth A&\ge\depth A/I_0\\
&\ge\min(\depth A/I_0 ,\depth (J\oplus J'))\\&=\dim A.
        \end{align*}
    Since the inequality $\dim A\ge \depth A$ always holds,  we get $\dim A=\depth A$ and $A$ is Cohen--Macaulay. We have the chain of inequalities
    \begin{align*}
        \dim A&=\depth A\bowtie^{f,g}(J,J')\\
        &=\min(\depth A/I_0 ,\depth (J\oplus J'))\\&\le\depth J\\&\le\dim A
    \end{align*}
    where the last inequality follows from the fact that $\depth J<\infty$. It must be the case that equality holds in every step of this chain and we get $\depth J=\dim A$ so $J$ is a maximal Cohen--Macaulay module. Similarly, $J'$ is a maximal Cohen--Macaulay module and $\depth A/I_0=\depth A$. Conversely, it is easy to see that $A\bowtie^{f,g}(J,J')$ is Cohen--Macaulay under the given conditions by Theorem \ref{dimdepthcalclemma}.\medskip

    (2) Let $A\bowtie^{f,g}(J,J')$ be Cohen--Macaulay. We will prove that $A/(\ker f\cap \ker g)$ is also Cohen--Macaulay. Using Theorem \ref{dimdepthcalclemma}, we obtain
        \begin{align*}
            \dim A/(\ker f\cap\ker g)&=\dim A\bowtie^{f,g}(J,J')\\&=\depth A\bowtie^{f,g}(J,J')\\
&=\min(\depth A/(\ker f\cap\ker g) ,\depth (J\oplus J'))\\&\le\depth A/(\ker f\cap\ker g)\\&\le\dim A/(\ker f\cap\ker g).
        \end{align*}
        Similarly, equality must hold in every step of this chain. Therefore we get that $A/(\ker f\cap\ker g)$ is Cohen--Macaulay. Let us assume further that $\depth J,\depth J'<\infty$. An analogous argument shows that $J$ and $J'$ are big Cohen--Macaulay $A/(\ker f\cap\ker g)$--modules. It is easy to see conversely that these assumptions imply $A\bowtie^{f,g}(J,J')$ is Cohen--Macaulay.
\end{proof}

\end{document}
